% sections/architecture.tex
% Included from main.tex via \input{sections/architecture}

\section{System Architecture}
\label{sec:architecture}

CARLA-Air integrates CARLA~\cite{carla} and AirSim~\cite{airsim} within a
single Unreal Engine~\cite{ue4} process through a minimal bridging layer that
resolves a fundamental engine-level initialization conflict while preserving
both platforms' native APIs, physics engines, and rendering pipelines intact.
Fig.~\ref{fig:arch_overview} presents the high-level runtime structure; the
following subsections elaborate each design decision.

\input{figures/fig_arch_overview}

\subsection{Plugin and Dependency Structure}
\label{sec:arch:plugins}

The system comprises two plugin modules that load sequentially during engine
startup.
The ground simulation plugin initializes first, establishing its world
management subsystems before any game logic executes.
The aerial simulation plugin declares a compile-time dependency on the ground
plugin, enabling \texttt{CARLAAirGameMode} to access the ground platform's
initialization interfaces during its own startup phase.
This dependency is strictly one-directional: no ground-platform source file
references any aerial component, preserving the upstream CARLA codebase's
update path without modification.
Two independent RPC servers run concurrently within the single
process---one per simulator---allowing the native Python clients of each
platform to connect without modification.

Version compatibility across the two upstream codebases, network configuration,
and port assignments are documented in Appendix~\ref{app:config}.

\subsection{The GameMode Conflict and Its Resolution}
\label{sec:arch:gamemode}

UE4 enforces a strict invariant: each world may have exactly one active game
mode.
CARLA's game mode orchestrates episode management, weather control, traffic
simulation, the actor lifecycle, and the RPC interface through a deep
inheritance chain.
AirSim's game mode performs a separate startup sequence---reading configuration
files, adjusting rendering settings, and spawning its flight actor.
Because the two game modes are unrelated by inheritance, assigning either to the
world map silently skips the other's initialization, rendering a large portion
of its API surface inoperative.

\paragraph{Architectural asymmetry.}
A structural difference between the two systems makes resolution tractable and
constitutes the central design insight of CARLA-Air.
CARLA's subsystems are tightly coupled to its game mode through inheritance and
privileged class relationships; they cannot be relocated outside the game mode
slot without invasive upstream refactoring.
AirSim's flight logic, by contrast, resides in a class derived from the
generic \emph{actor} base---not the game mode base---and can therefore be
spawned as a regular world actor at any point after world initialization.

\paragraph{Solution: single inheritance plus composition.}
We introduce \texttt{CARLAAirGameMode}, which inherits from CARLA's game mode
base and occupies the single available slot.
All ground simulation subsystems are thereby acquired through the standard UE4
lifecycle.
The aerial flight actor is then \emph{composed} into the world during the
engine's \textsc{BeginPlay} phase, after ground initialization is complete, and
never competes for the game mode slot.
Fig.~\ref{fig:gamemode_design} contrasts the naive conflict with this adopted
solution.

\input{figures/fig_gamemode_design}

\paragraph{Source footprint.}
The integration modifies exactly two files in the upstream CARLA source tree:
two previously private members are promoted to protected visibility and one
privileged class declaration is added.
All remaining integration code resides within the aerial plugin as purely
additive content.
The complete modification summary is provided in Appendix~\ref{app:source_mods}.

\subsection{Coordinate System Mapping}
\label{sec:arch:coords}

CARLA and AirSim employ incompatible spatial reference frames that must be
reconciled to co-register aerial and ground sensor data.
CARLA inherits UE4's left-handed system with X forward, Y right, and Z up, in
centimeters.
AirSim adopts a right-handed North-East-Down (NED) frame with X north, Y east,
and Z down, in meters.
Fig.~\ref{fig:coord_frames} illustrates both frames and their geometric
relationship.

\input{figures/fig_coord_frames}

Let $\mathbf{p} \in \mathbb{R}^3$ denote a point in the UE4 world frame and
$\mathbf{o}$ the shared world origin established during initialization.
The equivalent NED position is
%
\begin{equation}
  \mathbf{p}_{\mathrm{NED}}
  \;=\; \frac{1}{100}
  \begin{pmatrix} p_x - o_x \\ p_y - o_y \\ -(p_z - o_z) \end{pmatrix},
  \label{eq:pos_transform}
\end{equation}
%
\noindent where the scale factor converts centimeters to meters and the sign
reversal on the third component reflects the Z-axis inversion.
Because the X and Y axes are directionally aligned, no axis permutation is
required.

For orientation, let $q = (w, q_x, q_y, q_z)$ denote a unit quaternion in the
UE4 frame.
The equivalent NED quaternion is
%
\begin{equation}
  q_{\mathrm{NED}} \;=\; \bigl(w,\; q_x,\; q_y,\; -q_z\bigr),
  \label{eq:rot_transform}
\end{equation}
%
\noindent where negating $q_z$ accounts for the Z-axis reversal and the
associated change of frame handedness.
Eqs.~(\ref{eq:pos_transform}) and~(\ref{eq:rot_transform}) together fully
specify the pose transform, enabling consistent fusion of drone attitude from
the aerial API with vehicle heading from the ground API across all joint
simulation workflows.

\subsection{Asset Import Pipeline}
\label{sec:arch:assets}

CARLA-Air provides an extensible asset import pipeline that allows researchers
to bring custom robot platforms, UAV models, vehicles, and environment assets
into the shared simulation world. Imported assets are fully integrated into the
joint simulation environment: they participate in the same physics tick and
rendering pass as all built-in actors, respond to both ground and aerial API
calls, and are visible to all sensor modalities across both simulation
backends. This capability enables evaluation of custom hardware designs---such
as novel multirotor configurations or application-specific ground
robots---within realistic air-ground scenarios without modifying the core
CARLA-Air codebase. Fig.~\ref{fig:custom_assets} shows two examples of
user-imported assets operating within the platform.
